\section{Quantum Components}
\label{sec:quantum}

All circuits have at most 11 qubits and are simulated exactly as state
vectors, so they are classically tractable by construction. We do not ask
whether there is a quantum advantage. We ask whether a quantum inductive
bias helps class-incremental learning when compared with its classical twin.

\subsection{A Projected Quantum Kernel in Three Places}
\label{sec:pqk}

\paragraph{Kernel.}
A feature is mapped to $q=8$ rotation angles,
$\boldsymbol\theta(\mathbf z)=\pi\tanh\big(W\,\mathrm{LN}(\mathbf z)+\mathbf b\big)$.
The circuit loads the angles with $R_y$ rotations and then applies $L=2$
variational layers, each of trainable $R_y$ rotations followed by a ring of
CNOTs:
\begin{equation}
  |\psi(\mathbf z)\rangle=\prod_{l=1}^{L}\Big[U_{\mathrm{ent}}\textstyle\bigotimes_jR_y(\phi_{l,j})\Big]
  \textstyle\bigotimes_jR_y(\theta_j(\mathbf z))\,|0\rangle^{\otimes q}.
\end{equation}
Following~\cite{huang2021power}, we compare two states through their
one-qubit reduced density matrices $\rho_j(\mathbf z)$:
\begin{equation}
  k(\mathbf z,\mathbf z')=\exp\Big(-\gamma\sum_{j=1}^{q}\big\|\rho_j(\mathbf z)-\rho_j(\mathbf z')\big\|_F^2\Big).
  \label{eq:pqk}
\end{equation}
Local reduced matrices are used because global fidelities concentrate
exponentially in the number of qubits~\cite{thanasilp2024concentration}. The
bandwidth $\gamma$ is kept within a factor of 10 of the inverse median
distance.

\paragraph{Placements.}
\begin{enumerate}
  \item \emph{Steering metric.} In the base task, $k$ replaces the cosine
    similarity of the baseline's pairwise steering loss. The loss pulls
    same-class pairs towards $k=1$ and pushes other pairs below a margin.
    The kernel is trained with the adapter and then frozen.
  \item \emph{Separation metric.} In incremental tasks, the frozen $k$
    penalises similarity between new features and old class means. We
    evaluate it on either current or previous-model features.
  \item \emph{Classifier feature map.} Ridge heads are fitted on the
    classifier-alignment pseudo-features mapped through the PQK, and
    compared, within the same run, with the linear head and with classical
    maps on the same inputs (random Fourier features, random projections,
    QDA). A variant without CNOTs isolates the role of entanglement.
\end{enumerate}

\subsection{Class-Incremental Learning as State Discrimination}
\label{sec:qsd}

\paragraph{Encoding.}
Let $\mathbf m_0$ be the mean feature of task~0 and $s_0^2$ the trace of its
covariance; both are fixed after task~0. Each feature is lifted and
normalised:
\begin{equation}
  \mathbf v(\mathbf z)=\begin{bmatrix}(\mathbf z-\mathbf m_0)/s_0\\1\end{bmatrix}\in\mathbb R^{D},\qquad
  \psi(\mathbf z)=\frac{\mathbf v(\mathbf z)}{\|\mathbf v(\mathbf z)\|},\qquad D=769.
  \label{eq:encoding}
\end{equation}
$\psi$ is the amplitude vector of a 10-qubit pure state, zero-padded to
$2^{10}$ entries; the padding changes nothing below. The constant entry makes
the encoding injective: $\mathbf z$ can be recovered from $\psi$.

\paragraph{Class states from the memory.}
A class becomes the second-moment state of its lifted features:
\begin{equation}
  \rho_c=\frac{\mathbb E_c[\mathbf v\mathbf v^{\top}]}{\mathbb E_c\|\mathbf v\|^2},\qquad
  \mathbb E_c[\mathbf v\mathbf v^{\top}]=
  \begin{bmatrix}\tilde\Sigma_c+\tilde{\boldsymbol\mu}_c\tilde{\boldsymbol\mu}_c^{\top} & \tilde{\boldsymbol\mu}_c\\
  \tilde{\boldsymbol\mu}_c^{\top}&1\end{bmatrix},
  \label{eq:state}
\end{equation}
where $\tilde{\boldsymbol\mu}_c=(\boldsymbol\mu_c-\mathbf m_0)/s_0$ and
$\tilde\Sigma_c=\Sigma_c/s_0^2$. The state is computed from $\mathcal M_t$
alone. It therefore needs no stored features, takes no more memory than the
baseline, and follows the drift correction of the means. We also tested two
alternatives:
\begin{itemize}
  \item states built from the stored per-sample features, which is close
    to storing exemplars;
  \item those states transported between tasks by a channel
    $\rho\mapsto A_t\rho A_t^{\top}/\Tr(\cdot)$, with $A_t$ fitted in closed form
    on the current task's features before and after training.
\end{itemize}

\paragraph{Pretty-good measurement.}
With uniform priors, the PGM~\cite{hausladen1994pgm} uses the operators
\begin{equation}
  E_c=\frac1{C_t}\bar\rho_\delta^{-1/2}\rho_c\,\bar\rho_\delta^{-1/2},\qquad
  \bar\rho_\delta=\frac1{C_t}\sum_c\rho_c+\delta I,\qquad
  p(c\mid\mathbf z)=\frac{\psi^{\top}E_c\psi}{\sum_{c'}\psi^{\top}E_{c'}\psi},
  \label{eq:pgm}
\end{equation}
with $\delta=0.1/2^{10}$; $\psi^{\top}E_c\psi=\Tr(E_c\psi\psi^{\top})$ is the Born
probability of outcome $c$. Because CIL errors are mostly cross-task
(Section~\ref{sec:gap}), we also use a \emph{hierarchical} measurement:
\begin{itemize}
  \item it first measures the task, $E_T=\sum_{c\in\mathcal Y_T}E_c$;
  \item it then measures the class inside the task, whitened by the task
    state $\rho_T$.
\end{itemize}
The resulting score is $p(c\mid\mathbf z)=p(T\mid\mathbf z)\,p(c\mid\mathbf z,T)$
with $T=\tau(c)$.

\paragraph{Richer encoding.}
To give the measurement a non-linear feature space, we add a flag qubit and
a two-copy block:
$\big(|0\rangle|\psi\rangle+|1\rangle|\mathbf u\otimes\mathbf u\rangle\big)/\sqrt2$,
where $\mathbf u$ is a normalised random projection of $\mathbf z$ to 32
dimensions. This uses 11 qubits.

\subsection{What the Measurement Can and Cannot Do}
\label{sec:theory}

\begin{proposition}[valid states]\label{prop:state}
$\rho_c$ is positive semi-definite with unit trace. It is the mixture
$\mathbb E_c[\pi(\mathbf z)\psi\psi^{\top}]$ of the encoded pure states, with
weights $\pi(\mathbf z)=\|\mathbf v(\mathbf z)\|^2/\mathbb E_c\|\mathbf v\|^2$.
\end{proposition}
\begin{proof}
$\mathbf v\mathbf v^{\top}=\|\mathbf v\|^2\psi\psi^{\top}$. Taking expectations gives a
non-negative combination of rank-one projectors, with trace
$\mathbb E_c\|\mathbf v\|^2$.
\end{proof}

\begin{proposition}[unitary invariance]\label{prop:unitary}
For any unitary $U$ acting after the encoding ($\psi\mapsto U\psi$,
$\rho_c\mapsto U\rho_cU^{\dagger}$), every $p(c\mid\mathbf z)$ in
Eq.~\eqref{eq:pgm} is unchanged.
\end{proposition}
\begin{proof}
$\bar\rho_\delta\mapsto U\bar\rho_\delta U^{\dagger}$. Matrix functions commute with
unitary conjugation, so $E_c\mapsto UE_cU^{\dagger}$, and
$(U\psi)^{\dagger}UE_cU^{\dagger}U\psi=\psi^{\dagger}E_c\psi$.
\end{proof}

\begin{proposition}[shrinkage is inert]\label{prop:shrink}
Replacing every $\rho_c$ by $(1-\varepsilon)\rho_c+\varepsilon\bar\rho$, with
$\varepsilon\in[0,1)$, leaves $\arg\max_cp(c\mid\mathbf z)$ unchanged.
\end{proposition}
\begin{proof}
$\bar\rho$ and $\bar\rho_\delta$ do not change. Each numerator becomes
$(1-\varepsilon)\psi^{\top}E_c\psi+\varepsilon\psi^{\top}\bar E\psi$, with
$\bar E=\frac1{C_t}\bar\rho_\delta^{-1/2}\bar\rho\,\bar\rho_\delta^{-1/2}$. The
denominator is unchanged. This is one increasing affine map, applied to
every class alike.
\end{proof}

\paragraph{Consequences.}
\begin{itemize}
  \item A variational circuit placed after the amplitude encoding is a
    unitary. By Proposition~\ref{prop:unitary}, training one cannot change
    the classifier.
  \item By Proposition~\ref{prop:shrink}, the usual way of regularising
    class states has no effect, so $\delta$ is the only effective
    regulariser.
  \item The PGM coincides with the least-squares measurement~\cite{eldar2001srm}.
    On second moments it is therefore the quantum counterpart of
    least-squares classification, which gives its classical twin below.
\end{itemize}
What remains as a quantum degree of freedom is the encoding itself.

\subsection{Measurement Fusion and its Classical Twin}
\label{sec:fusion}

\paragraph{Classical twin.}
Treat each class as $\mathcal N(\boldsymbol\mu_c,\Sigma_c)$ and regress one-hot
targets with squared loss and ridge penalty. The normal equations give
\begin{equation}
  \Big(\sum_c\big(\Sigma_c+\boldsymbol\mu_c\boldsymbol\mu_c^{\top}\big)+\lambda I\Big)W=
  [\boldsymbol\mu_1,\dots,\boldsymbol\mu_{C_t}],\qquad g^{\mathrm R}_c(\mathbf z)=\mathbf z^{\top}W_{:,c},
  \label{eq:ridge}
\end{equation}
with $\lambda=0.1\,\mathrm{tr}(\cdot)/d$. Like the PGM, this uses only
$\mathcal M_t$. Closed-form heads of this kind are known in pre-trained
CIL~\cite{mcdonnell2023ranpac}.

\paragraph{Fusion.}
A generative score $g$ (the PGM log-probability $g^{\mathrm Q}=\log p$, or
$g^{\mathrm R}$) is added to the linear-head logits:
\begin{equation}
  \hat y(\mathbf z)=\arg\max_c\;\ell_c(\mathbf z)+w\,\kappa\,g_c(\mathbf z),\qquad
  \kappa=\frac{\frac1n\sum_i\mathrm{sd}_c\,\ell_c(\mathbf x_i)}{\frac1n\sum_i\mathrm{sd}_c\,g_c(\mathbf x_i)}.
  \label{eq:fusion}
\end{equation}
Here $\mathrm{sd}_c$ is the spread across classes, and $\mathbf x_i$ are the
training images of the current task. $\kappa$ needs no labels. The weight
$w\in\{0.25,0.5,1,2,4\}$ is chosen once on the tuning group and frozen.

For the PGM, the softmax of the fused score is
$p_{\mathrm{lin}}(c\mid\mathbf z)\,p(c\mid\mathbf z)^{w\kappa}$ up to normalisation. For
an integer $k$, $p(c\mid\mathbf z)^k$ is the Born probability of the product
measurement $E_c^{\otimes k}$ on $k$ copies of the state. The weight
therefore acts as an effective number of measured copies.

\begin{proposition}[scale invariance]\label{prop:scale}
For any $a>0$ and any class-independent $b(\mathbf z)$, replacing $g_c$ by
$a\,g_c+b$ leaves $\hat y$ in Eq.~\eqref{eq:fusion} unchanged.
\end{proposition}
\begin{proof}
$\kappa\mapsto\kappa/a$, so $\kappa g_c\mapsto\kappa g_c+\kappa b/a$, and the
added term does not depend on $c$.
\end{proof}
\noindent Because of Proposition~\ref{prop:scale}, the same $w$ means the same
thing for both sources and on every dataset. This is what makes a frozen
weight a fair comparison between the PGM and its twin.

\paragraph{Cost.}
Neither source stores anything beyond $\mathcal M_t$.
\begin{itemize}
  \item Ridge: one $d\times d$ solve per task.
  \item PGM: one $D\times D$ eigendecomposition per task, then one quadratic
    form per class and test image. At the last ImageNet-R task this is
    about $1.4$ TFLOP, under one second on a GPU.
\end{itemize}
